\documentclass[11pt,a4paper]{article}
\usepackage[margin=27mm]{geometry}
\usepackage{amsmath,amssymb,amsthm,fontspec,unicode-math,hyperref}
\setmainfont{STIXGeneral}\setmathfont{STIX Two Math}
\setlength{\emergencystretch}{2em}
\newtheorem{theorem}{Theorem}\newtheorem{proposition}[theorem]{Proposition}
\title{Chiral polyhedra and one projective reversal}
\author{The human\_compression project}\date{11 October 2026}
\begin{document}\maketitle
\begin{abstract}
We recover the complete principal theorem of Elliott, Leemans and
O'Reilly-Regueiro: neither $\mathrm{PSL}_2(q)$ nor $\mathrm{PGL}_2(q)$
is the full automorphism group of a chiral polyhedron. The common
intertwiner theorem gives one reversal for every generating pair,
including characteristic two and the small fields. We also give the
precise characteristic qualification for the source's auxiliary count.
The full common theory is an explicit dependency of this application.
\end{abstract}

\begin{theorem}
For every prime power $q$, no abstract chiral polyhedron has full
automorphism group $\mathrm{PSL}_2(q)$ or $\mathrm{PGL}_2(q)$.
\end{theorem}
\begin{proof}
Let $a,b$ be its distinguished rotations, with $(ab)^2=1$.
Choose invertible matrix representatives for $a$ and $ab$.
Their generated projective group contains the opposite elementary
unipotents in $\mathrm{PSL}_2(q)$, whose unique algebraic fixed points
differ. Thus the representatives have no common invariant line even
over an algebraic closure. By the simultaneous-reversal theorem of
[U], their additive commutator $X$ is invertible, $X^2$ is scalar,
and conjugation by $[X]$ inverts both projectivities. It therefore
inverts $a$ and fixes $ab$, sending $b$ to $a^2b$.
This involutory automorphism normalizes both projective groups.
It is exactly the reflection forbidden by the retained rotation-group
criterion, a contradiction. The argument includes $q=2,3$; it does
not replace these fields by an unproved computer lookup.
\end{proof}

The matrices underlying this theorem are fully explicit:
$X=AB-BA$ satisfies $XA=\operatorname{adj}(A)X$ and
$XB=\operatorname{adj}(B)X$. The common proof establishes both these
identities and the nonzero determinant from irreducibility.
No finite classification of generating pairs is needed. This replaces
the source's three separate fixed-point cases (Lemmas~4.4--4.6) by a
single theorem with strictly weaker hypotheses: the product relation
and cyclic intersection are not needed to construct the simultaneous
inverter, only to interpret it as a polyhedral reflection.

\section{The point stabilizer and the auxiliary correction}
We record the elementary projective-line inputs rather than retain an
unstated exception. A nonidentity element whose order divides $q=p^d$
is unipotent and has order $p$, by its two-dimensional Jordan form.
All such elements are conjugate in $\mathrm{PGL}_2(q)$ and lie in
$\mathrm{PSL}_2(q)$. A noninvolutory semisimple projectivity has two
rational fixed points in the split case and none in the nonsplit case;
its eigenvalue ratio has order dividing $q-1$ or $q+1$, respectively.
Conversely these divisibilities distinguish the cases since
$\gcd(q-1,q+1)\le2$. A projectivity interchanging two distinct points
is antidiagonal after sending the points to zero and infinity, and its
square is scalar. These are the precise forms of source Lemma~4.2;
the identity must be excluded from its first assertion.

Normalize a nonidentity unipotent to $u:x\mapsto x+1$.
An inverting projectivity fixes its unique fixed point infinity and
therefore has form $x\mapsto ax+b$. The conjugation equation forces
$a=-1$. In odd characteristic every $b\in F_q$ gives a nonidentity
involution, so the source Lemma~4.3 count is $q$. In characteristic two
these maps are translations: $b=0$ gives the identity and the remaining
$q-1$ give involutions. Thus the literal unrestricted count is $q-1$
in characteristic two. This correction changes neither the principal
theorem nor the uniform reversal proof above.

\section{Source and evidence boundary}
The original result and its geometric approach are due to
Elliott--Leemans--O'Reilly-Regueiro [ELO]. We bind their complete
eight-page authored edition. Its Table~1 is explicitly borrowed from
Hartley--Hubard--Leemans' atlas; that motivating census is not claimed
as a new calculation here. Its Conjecture~3.1 is the principal theorem;
its subsequent rank and rank-four-existence questions were addressed in
Leemans--Moerenhout--O'Reilly-Regueiro, the companion application in this
same topic. They are recorded historically, not counted as new answers.

The common paper [U] and the retained group--polytope correspondence
are the full proof dependencies. The new Lean module certifies actual
commutator intertwining and scalar-square identities. Absolute
irreducibility and the projective and geometric interpretation remain
written proofs; no complete Lean polyhedron theorem is claimed.
There is no required finite program for this particular principal result.

\begin{thebibliography}{3}
\bibitem[U]{U} The human\_compression project,
\emph{Common theory for projective markings and rotations}, current complete common paper.
\bibitem[ELO]{ELO} D. Elliott, D. Leemans and E. O'Reilly-Regueiro,
\emph{Chiral polyhedra and projective lines}, Int. J. Math. Stat.
16(1) (2015), 45--52; bound eight-page authored edition.
\bibitem{SW} E. Schulte and A. I. Weiss, \emph{Chiral polytopes},
Applied Geometry and Discrete Mathematics, DIMACS 4 (1991), 493--516;
retained rotation correspondence.
\end{thebibliography}
\end{document}
